\color{unimain}{\subsection{Максимальная токовая защита компенсационной обмотки}}
\color{black}

\begin{enumerate}

\item
Функция выполнена двухступенчатой и позволяет реализовать защиту:
\begin{itemize}
\item от междуфазных КЗ в КОР и на ее выводах (включая отходящие присоединения ошиновки НН), чувствительную к межвитковым замыканиям в КОР, с подключением измерительных цепей защиты на вторичные фазные токи ТТ, установленных непосредственно в фазах КОР;
\item от междуфазных КЗ в КОР и на ее выводах (включая отходящие присоединения ошиновки НН), с подключением измерительных цепей защиты на разность вторичных фазных токов ТТ, установленных непосредственно в фазах КОР:
\end{itemize}
\begin{equation}
i'_{A} = \frac{i_{A}-i_{B}}{\sqrt3}; \; i'_{B} = \frac{i_{B}-i_{C}}{\sqrt3}; \; i'_{C} = \frac{i_{C}-i_{A}}{\sqrt3}
\end{equation}

\item
Тип подключения защиты определяется программной накладкой \textbf{«ИО МТЗ КО»}. Защита имеет регулируемые значения по току и времени срабатывания. Каждая из ступеней защиты может выполнена с блокировкой по второй гармонике (отношению уровня второй гармоники к первой). Данная блокировка реализована на базе модуля \textbf{«ОВ БНТ»}. Предусмотрена блокировка по току 3I0 стороны ВН реактора.

\item
Алгоритм работы МТЗ КО-1 приведен на рисунке \ref{pict_mtzko:MTZ_KO}, алгоритм работы МТЗ КО-2 выполнен аналогично. Уставки МТЗ КО приведены в таблице \ref{app_set:mtzko}.

\medskip
\begin{figure}[H]
\centering
\includegraphics[width=1\textwidth,height=1\textheight,keepaspectratio]{\fbpath/11.02 УШР/МТЗ КО/img/MTZ_KO.pdf}
\captionof{figure}{Функциональная схема логики МТЗ КО-1}\label{pict_mtzko:MTZ_KO}
\end{figure}


\end{enumerate}